% ============================================================
%  MCM/ICM paper skeleton  (compiles with pdfLaTeX, the Overleaf default)
%  Files needed in the same folder:  main.tex  easymcm.sty
%  Put figures in a folder named  figure/  (PNG or PDF, no spaces in file names)
%  Keep this file in English only: Chinese characters will break pdfLaTeX.
% ============================================================
\documentclass[12pt]{article}
\usepackage[0000000]{easymcm}     % <- replace 0000000 with your team control number
\usepackage{mathptmx}             % Times-like font
\usepackage{enumitem}
\problem{A}                       % <- problem chosen: A-F
\title{Your Paper Title: Specific and Informative}

\begin{document}

% ---------------- Summary Sheet (page 1) ----------------
\begin{abstract}
One or two sentences of background and task.

For Problem 1, we build ... because ... The result shows that ... (give numbers).

For Problem 2, ...

Sensitivity analysis shows that ... The model remains reliable when ...

Our key finding / recommendation is ...

\vspace{5pt}
\textbf{Keywords}: keyword1; keyword2; keyword3; keyword4
\end{abstract}
\maketitle

% ---------------- Contents ----------------
\tableofcontents

% ---------------- Body ----------------
\section{Introduction}
\subsection{Background}
\subsection{Restatement of the Problem}
\subsection{Our Work}
% Put the overall flowchart of the paper here.

\section{Assumptions and Justifications}
\begin{enumerate}[label=\textbf{Assumption \arabic*:}, leftmargin=*, itemsep=4pt]
    \item \textbf{State the assumption.}\\
          \textit{Justification:} why it is reasonable.
    \item \textbf{State the assumption.}\\
          \textit{Justification:} why it is reasonable.
\end{enumerate}

\section{Notations}
\begin{table}[H]
    \centering
    \caption{Notations}\label{tab:notation}
    \begin{tabular}{cl}
        \toprule
        Symbol & Description \\
        \midrule
        $r$ & Intrinsic growth rate \\
        $K$ & Carrying capacity \\
        $x(t)$ & Population at time $t$ \\
        \bottomrule
    \end{tabular}
\end{table}

\section{Model I: Descriptive Title}
\subsection{Model Development}
Start from the simplest version, then refine. A numbered equation:
\begin{equation}\label{eq:logistic}
    \frac{\mathrm{d}x}{\mathrm{d}t} = r x \left(1 - \frac{x}{K}\right)
\end{equation}
Equation~\eqref{eq:logistic} is referenced like this. Symbols are listed in Table~\ref{tab:notation}.

\subsection{Solution and Results}
A figure (replace the placeholder box with \verb|\includegraphics|):
\begin{figure}[H]
    \centering
    % \includegraphics[width=0.75\textwidth]{fig_ode_fit}   % file: figure/fig_ode_fit.png
    \fbox{\rule{0pt}{4cm}\rule{0.6\textwidth}{0pt}}
    \caption{What the figure shows, in one sentence.}\label{fig:demo}
\end{figure}
Figure~\ref{fig:demo} indicates that ...

\subsection{Validation}

\section{Model II: Descriptive Title}

\section{Sensitivity Analysis}

\section{Strengths and Weaknesses}
\subsection{Strengths}
\subsection{Weaknesses}

\section{Conclusion}

% ---------------- Letter / Memo (only if the problem asks for one) ----------------
\begin{letter}{Letter to the Decision Maker}
\noindent Dear Sir or Madam,

Plain language, no formulas, concrete recommendations.

\vspace{1em}
\noindent Sincerely,\\
Team \control
\end{letter}

% ---------------- References ----------------
\clearpage
\begin{thebibliography}{99}
    \bibitem{ref1} Author, A. (2024). Title of the paper. \textit{Journal Name}, 12(3), 45--67.
    \bibitem{ref2} Organization. (2025). Dataset name. \url{https://example.org}. Accessed 30 Jan 2027.
\end{thebibliography}

% ---------------- Appendix (counts toward the 25 pages: keep it short) ----------------
\begin{subappendices}
\section{Appendix: Core Code}
\begin{lstlisting}[language=Python]
import numpy as np
print(np.pi)
\end{lstlisting}
\end{subappendices}

% ---------------- Report on Use of AI (after the 25 pages, not counted) ----------------
\clearpage
\section*{Report on Use of AI}
\begin{enumerate}[leftmargin=*]
    \item \textbf{Tool:} name and version of the model.\\
          \textbf{Query:} what we asked.\\
          \textbf{Output and how we used it:} what it returned, what we kept, how we verified it.
\end{enumerate}

\end{document}
